%LTeX: language=fr-FR
\documentclass[12pt]{article}
\usepackage[fontsize=13pt]{scrextend}
\usepackage{amsmath,amssymb}
\usepackage[french]{babel}
\usepackage{enumitem}
\usepackage{currfile}
\usepackage{xcolor}
\usepackage{tasks}
\usepackage{pifont}
\usepackage{tikz}
\usetikzlibrary{arrows.meta,patterns}
\usepackage{geometry}
\geometry{left=1.5cm, right=1.5cm, top=1cm, bottom=1cm}
\setlength{\parindent}{0cm}
\setlength{\parskip}{2pt}
\pagestyle{empty}
\input{\currfiledir ../def.tex}
\def \Chapitre {Nombres et fonctions}
\def \Document {Contrôle}


\usepackage{tcolorbox}
\tcbuselibrary{skins}
\definecolor{bandeau}{HTML}{E8EEF7}
% Bandeau de titre de chaque exercice (même style que les contrôles de 1re ST2S)
\newtcbox{\exobox}{colback=bandeau, colframe=bandeau, boxrule=0pt, arc=0pt, left=4pt, right=4pt, top=3pt, bottom=3pt, fontupper=\bfseries}
\newcounter{exo}
% Espace avant chaque exercice (élastique dans les sujets, pour répartir les exercices sur la page)
\newcommand{\exoskip}{\vspace{1.4em}}
% \begin{exercice}{titre}{barème}
\newenvironment{exercice}[2]{%
  \stepcounter{exo}\par\ifnum\value{exo}=1 \vspace{0.8em}\else\exoskip\fi
  \noindent\exobox{Exercice \theexo{} --- #1 (#2 points)}\par\vspace{0.3em}
}{\par}
\setlist[enumerate]{label=\arabic*., itemsep=2pt, topsep=3pt}
\colorlet{sol}{blue!70!black}
\newcommand{\sol}[1]{\textcolor{sol}{#1}}
\newcommand{\bareme}[1]{\par\vspace{6pt}\noindent\textcolor{sol}{\small #1}}

\begin{document}
\renewcommand{\exoskip}{\vspace{1.4em plus 1fill}}
\setlist[enumerate]{label=\arabic*., itemsep=9pt, topsep=6pt}
\begin{center}
  {\large\bfseries Contrôle (B) --- Nombres et fonctions}\\[0.3em]
  {\itshape Seconde \quad $\bullet$ \quad Durée~: 55 minutes}
\end{center}
\noindent\textit{Toutes les réponses doivent être justifiées.}
\begin{exercice}{Automatismes}{5}
Les questions sont indépendantes. Détailler les calculs.
\begin{enumerate}
  \item Calculer $\dfrac{6}{25}\times\dfrac{10}{9}$. Donner le résultat sous forme irréductible.
  \item Résoudre l'inéquation $-2x+10<0$. Donner les solutions sous forme d'intervalle.
  \item Calculer $75\,\%$ de $160$.
  \item Justifier chaque appartenance.
  \begin{tasks}[label=\alph*), label-width=1.4em, item-indent=1.8em](4)
    \task $\dfrac{-21}{7}\in\mathbb{Z}$
    \task $\dfrac{9}{20}\in\mathbb{D}$
    \task $-1{,}8\in\mathbb{Q}$
    \task $\sqrt{64}\in\mathbb{N}$
  \end{tasks}
\end{enumerate}

\end{exercice}

\begin{exercice}{Lecture graphique}{5}
La courbe $\mathcal{C}_f$ ci-contre représente une fonction $f$. Répondre par lecture graphique.
\par\smallskip
\begin{minipage}[t]{0.46\linewidth}\raggedright\hyphenpenalty=10000\exhyphenpenalty=10000
\begin{enumerate}
  \item Quel est l'ensemble de définition de $f$~?
  \item Donner $f(-1)$ et l'image de $5$ par $f$.
  \item Déterminer les antécédents de $0$ par $f$.
  \item Résoudre l'inéquation $f(x) \geqslant 0$.
\end{enumerate}
\end{minipage}\hfill
\begin{minipage}[t]{0.52\linewidth}
\vspace{0pt}\centering
\begin{tikzpicture}[x=0.98cm, y=0.9cm, lab/.style={fill=white, fill opacity=0.75, text opacity=1, inner sep=0.6pt}]
  \draw[gray!55, thin, xstep=1, ystep=1] (-4,-4) grid (5,2);
  \draw[-{Latex[length=2mm]}, thick] (-4,0) -- (5+0.5,0);
  \draw[-{Latex[length=2mm]}, thick] (0,-4) -- (0,2+0.5);
  \draw[thick, blue!60!black] plot coordinates {(-4.00,1.000) (-3.95,0.950) (-3.90,0.902) (-3.85,0.853) (-3.80,0.806) (-3.75,0.758) (-3.70,0.711) (-3.65,0.664) (-3.60,0.617) (-3.55,0.569) (-3.50,0.522) (-3.45,0.474) (-3.40,0.425) (-3.35,0.376) (-3.30,0.326) (-3.25,0.274) (-3.20,0.222) (-3.15,0.169) (-3.10,0.114) (-3.05,0.058) (-3.00,0.000) (-2.95,-0.061) (-2.90,-0.128) (-2.85,-0.199) (-2.80,-0.274) (-2.75,-0.354) (-2.70,-0.437) (-2.65,-0.523) (-2.60,-0.613) (-2.55,-0.705) (-2.50,-0.799) (-2.45,-0.895) (-2.40,-0.993) (-2.35,-1.092) (-2.30,-1.192) (-2.25,-1.293) (-2.20,-1.394) (-2.15,-1.495) (-2.10,-1.595) (-2.05,-1.695) (-2.00,-1.793) (-1.95,-1.891) (-1.90,-1.986) (-1.85,-2.079) (-1.80,-2.169) (-1.75,-2.257) (-1.70,-2.342) (-1.65,-2.423) (-1.60,-2.500) (-1.55,-2.573) (-1.50,-2.641) (-1.45,-2.705) (-1.40,-2.763) (-1.35,-2.816) (-1.30,-2.863) (-1.25,-2.903) (-1.20,-2.937) (-1.15,-2.964) (-1.10,-2.984) (-1.05,-2.996) (-1.00,-3.000) (-0.95,-2.996) (-0.90,-2.984) (-0.85,-2.964) (-0.80,-2.938) (-0.75,-2.904) (-0.70,-2.864) (-0.65,-2.817) (-0.60,-2.765) (-0.55,-2.707) (-0.50,-2.644) (-0.45,-2.576) (-0.40,-2.503) (-0.35,-2.426) (-0.30,-2.346) (-0.25,-2.262) (-0.20,-2.174) (-0.15,-2.084) (-0.10,-1.992) (-0.05,-1.897) (0.00,-1.800) (0.05,-1.702) (0.10,-1.602) (0.15,-1.502) (0.20,-1.402) (0.25,-1.301) (0.30,-1.200) (0.35,-1.100) (0.40,-1.001) (0.45,-0.903) (0.50,-0.806) (0.55,-0.712) (0.60,-0.619) (0.65,-0.529) (0.70,-0.442) (0.75,-0.359) (0.80,-0.278) (0.85,-0.202) (0.90,-0.130) (0.95,-0.063) (1.00,0.000) (1.05,0.061) (1.10,0.123) (1.15,0.186) (1.20,0.250) (1.25,0.316) (1.30,0.382) (1.35,0.448) (1.40,0.515) (1.45,0.583) (1.50,0.650) (1.55,0.717) (1.60,0.785) (1.65,0.852) (1.70,0.918) (1.75,0.984) (1.80,1.050) (1.85,1.114) (1.90,1.177) (1.95,1.239) (2.00,1.300) (2.05,1.359) (2.10,1.417) (2.15,1.473) (2.20,1.526) (2.25,1.578) (2.30,1.628) (2.35,1.675) (2.40,1.719) (2.45,1.761) (2.50,1.800) (2.55,1.836) (2.60,1.869) (2.65,1.898) (2.70,1.924) (2.75,1.947) (2.80,1.966) (2.85,1.980) (2.90,1.991) (2.95,1.998) (3.00,2.000) (3.05,1.992) (3.10,1.968) (3.15,1.929) (3.20,1.877) (3.25,1.812) (3.30,1.736) (3.35,1.649) (3.40,1.552) (3.45,1.446) (3.50,1.333) (3.55,1.213) (3.60,1.088) (3.65,0.958) (3.70,0.824) (3.75,0.688) (3.80,0.549) (3.85,0.411) (3.90,0.272) (3.95,0.135) (4.00,0.000) (4.05,-0.140) (4.10,-0.292) (4.15,-0.456) (4.20,-0.629) (4.25,-0.812) (4.30,-1.004) (4.35,-1.203) (4.40,-1.408) (4.45,-1.619) (4.50,-1.833) (4.55,-2.052) (4.60,-2.272) (4.65,-2.494) (4.70,-2.716) (4.75,-2.938) (4.80,-3.157) (4.85,-3.374) (4.90,-3.588) (4.95,-3.797) (5.00,-4.000)};
  \fill[blue!60!black] (-4, 1) circle (2.2pt) (5, -4) circle (2.2pt);
  \foreach \x in {-4,...,-1,1,2,...,5} {\draw[thick] (\x,0.08) -- (\x,-0.08); \node[below, lab, yshift=-1.5pt] at (\x,0) {\footnotesize $\x$};}
  \foreach \y in {-4,...,-1,1,2,...,2} {\draw[thick] (0.08,\y) -- (-0.08,\y); \node[left, lab, xshift=-1.5pt] at (0,\y) {\footnotesize $\y$};}
  \node[below left, lab, xshift=-1pt, yshift=-1pt] at (0,0) {\footnotesize $0$};
  \node[lab, text=blue!60!black] at (2.2, 2.6) {$\mathcal{C}_f$};
\end{tikzpicture}
\end{minipage}

\end{exercice}

\vfill
\noindent\hfill\textbf{Tournez la page}\enspace{\Large\ding{228}}
\newpage
\begin{exercice}{Problème}{6}
Un loueur de trottinettes propose deux formules :
\begin{itemize}[label=\textbullet, itemsep=1pt, topsep=2pt]
  \item \textbf{Formule A} : $5$~€ par heure ;
  \item \textbf{Formule B} : $12$~€ de frais de dossier, puis $3$~€ par heure.
\end{itemize}
On note $x$ la durée de location en heures, avec $0 \leqslant x \leqslant 12$. Le prix payé est $f(x) = 5x$ avec la formule A, et $g(x) = 12 + 3x$ avec la formule B.

\begin{minipage}[t]{0.5\linewidth}\raggedright\hyphenpenalty=10000\exhyphenpenalty=10000
\begin{enumerate}
  \item Calculer $f(4)$ et $g(4)$. Quelle formule est la moins chère pour $4$ heures~?
  \item Le point $M(10\,;\,40)$ appartient-il à la courbe $\mathcal{C}_g$~? Justifier par un calcul.
  \item Les courbes de $f$ et $g$ sont tracées ci-contre. Résoudre graphiquement $f(x) = g(x)$. Interpréter le résultat.
  \item Résoudre graphiquement $f(x) < g(x)$. Interpréter le résultat.
\end{enumerate}
\end{minipage}\hfill
\begin{minipage}[t]{0.48\linewidth}
\vspace{0pt}\centering
\begin{tikzpicture}[x=0.52cm, y=0.09cm]
  \draw[gray!35, very thin, xstep=1, ystep=5] (0,0) grid (12,60);
  \draw[-{Latex[length=2mm]}] (0,0) -- (12+0.6,0) node[below=9pt, anchor=north east] {\scriptsize heures};
  \draw[-{Latex[length=2mm]}] (0,0) -- (0,63) node[above] {\scriptsize prix (€)};
  \foreach \x in {2,4,...,12} \node[below] at (\x,0) {\scriptsize $\x$};
  \foreach \y in {10,20,...,60} \node[left] at (0,\y) {\scriptsize $\y$};
  \node[below left] at (0,0) {\scriptsize $0$};
  \begin{scope}
    \clip (0,0) rectangle (12,60);
    \draw[very thick] plot[domain=0:12, samples=2] (\x, {5*\x});
    \draw[very thick, dashed] plot[domain=0:12, samples=2] (\x, {12+3*\x});
  \end{scope}
  \node at (10.6, 57) {$\mathcal{C}_f$};
  \node at (11.5, 39.5) {$\mathcal{C}_g$};
\end{tikzpicture}
\end{minipage}

\end{exercice}

\begin{exercice}{Fonction et géométrie}{4}
\begin{minipage}[t]{0.6\linewidth}
\vspace{0pt}\raggedright\hyphenpenalty=10000\exhyphenpenalty=10000
$ABCD$ est un carré de côté $8$~cm.

$M$ est un point du segment $[AB]$.

On note $x$ la longueur $AM$, en cm.

$AMNP$ est un carré, avec $P$ sur le segment $[AD]$.

On note $\mathcal{A}(x)$ l'aire de la partie hachurée, en cm$^2$.
\begin{enumerate}
  \item Dans quel intervalle le nombre $x$ peut-il varier~?
  \item Justifier que $\mathcal{A}(x) = 64 - x^2$.
  \item Calculer $\mathcal{A}(3)$. Interpréter le résultat.
  \item Montrer que $5$ est un antécédent de $39$ par $\mathcal{A}$. Interpréter le résultat.
\end{enumerate}
\end{minipage}\hfill
\begin{minipage}[t]{0.36\linewidth}
\vspace{0pt}\centering
\begin{tikzpicture}[scale=0.47]
  \fill[pattern=north east lines] (0,0) rectangle (8,8);
  \fill[white] (0,0) rectangle (3,3);
  \draw[very thick] (0,0) -- (8,0) -- (8,8) -- (0,8) -- cycle;
  \draw[thick] (3,0) -- (3,3) -- (0,3);
  \node[below left] at (0,0) {$A$};
  \node[below right] at (8,0) {$B$};
  \node[above right] at (8,8) {$C$};
  \node[above left] at (0,8) {$D$};
  \node[below right] at (3,0) {$M$};
  \node[left] at (0,3) {$P$};
  \node[fill=white, inner sep=1pt, above right] at (3,3) {$N$};
  \draw[densely dashed, gray] (0,-0.2) -- (0,-1.7) (3,-0.2) -- (3,-1.7);
  \draw[{Bar[width=2.5mm]Latex[length=2mm]}-{Latex[length=2mm]Bar[width=2.5mm]}, thick] (0,-1.5) -- node[below=1pt] {$x$} (3,-1.5);
\end{tikzpicture}
\end{minipage}

\end{exercice}
\vfill
\end{document}
